\section{多会话研报流水线}
\label{sec:pipeline}

\begin{figure}[t]
  \centering
  \figslot{fig5_pipeline}{6.0cm}
  \caption{研报流水线。确定性编排器在仓库目录树之外为单家公司搭建隔离工作区，依次启动五个互不继承上下文的会话：研究助理（阶段 1）→ 分析师（阶段 2～7）→ 主笔（阶段 8）→ 复核编辑（阶段 9）→ 排版设计师（阶段 10～11）。会话之间只通过工作区文件交接；每个阶段结束都运行 \code{status}，由 11 道程序闸门判定能否进入下一阶段。}
  \label{fig:pipeline}
\end{figure}

前几节的引擎回答了“数字怎么算对、怎么不被写错”。本节回答另一个问题：怎样让模型写出一份有深度、像分析师写的研报。我们的答案是把研报生产拆成五个角色、十一个阶段，用确定性的编排器和程序闸门把它们串起来（图 \ref{fig:pipeline}）。

\subsection{隔离工作区}
\code{report\_pipeline build} 为一家公司在仓库目录树之外（默认 \file{\textasciitilde/finagent-samples/<公司>\_<年度>}）建立工作区，只放入四样东西：三份技能说明、计算引擎（\textbf{不含}评测模块 \file{finagent/evaluation/}）、这家公司的年报，以及人提供的外部输入（例如参考股价及其来源）。

隔离的理由来自实际教训：仓库里有其他公司的抽取文件、分析框架、文稿和交付物，模型在仓库里工作时会照抄别家公司的框架，并浪费大量上下文去读无关文件；同时仓库里有评分标准，放在生成方能读到的地方就违背了决策 D3。工作区的权限配置因此做了两道围栏：\textbf{禁止读取}仓库所在目录及其上级目录；\textbf{禁止修改}引擎、技能说明、年报原件与工作区说明。允许的工具调用以白名单形式在启动无头会话时传入，包括以 \code{./py -m finagent} 开头的引擎命令、对 \file{data/} 的读写、联网检索与抓取等。

\subsection{五个角色}
表 \ref{tab:roles} 列出五个会话及其默认模型。每个会话由一段简短的角色提示启动（附录 \ref{app:prompts} 给出要点），具体做法写在技能说明和工作区说明（\file{CLAUDE.md} / \file{AGENTS.md}）里。会话之间不继承上下文；每个会话完成自己的阶段后运行 \code{status}，闸门通过即停止并汇报，由编排器决定是否启动下一个会话。已经通过的阶段在重跑时自动跳过。

\begin{table}[t]
\centering
\caption{五个生产会话。生成用 Opus 5.5，复核用 Sonnet 5.5；均可通过环境变量替换，执行器也可换成 Codex。}
\label{tab:roles}
\small
\begin{tabularx}{\linewidth}{@{}p{1.7cm}p{1.2cm}p{1.4cm}X@{}}
\toprule
会话 & 模型 & 阶段 & 职责与边界 \\
\midrule
研究助理 & Opus & 1 & 联网搜集年报之后的信息，登记 \file{sources.json}、写资料包 \file{dossier.md}；不做财务抽取、不写分析框架 \\
分析师 & Opus & 2～7 & 财务抽取与校验、分析框架（含深度专题）、深度调研、经营数据、估值假设、估值运行；不写文稿 \\
主笔 & Opus & 8 & 读框架、资料包、专题笔记与运行产物，写“年报点评 + 深度研究”文稿，过护栏与篇幅闸门 \\
复核编辑 & Sonnet & 9 & 没有参与写作；只改措辞与结构，过数字锁；拿不准但重要的事项写入人工复核清单 \\
排版设计师 & Opus & 10～11 & 文稿已定稿，不改正文；写本公司的排版脚本、渲染 PDF、逐页检查，写交付说明 \\
\bottomrule
\end{tabularx}
\end{table}

\subsection{十一道闸门}
\label{sec:gates}
闸门由 \code{status} 统一判定，全部是程序检查（表 \ref{tab:gates}）。三条横向规则使闸门难以被“绕过”：
\begin{enumerate}
  \item \textbf{新鲜度。}校验结果必须晚于被校验文件的最后修改时间：例如 \file{facts.json} 改过而未重新校验，阶段 2 判未通过；估值运行早于任何输入文件的修改，阶段 7 判为“已过期”；PDF 早于文稿或排版脚本，阶段 10 判未通过。
  \item \textbf{核对的是当前稿。}阶段 8 要求护栏结果对应的文稿与当前文稿逐字相同，并且护栏带上了外部来源登记。
  \item \textbf{前序阻断。}只要前序阶段未通过，后续阶段即使文件齐全也一律判为未通过，避免拿过期产物当成果。
\end{enumerate}

\begin{table}[t]
\centering
\caption{十一道程序闸门及其阈值（常量定义于 \file{report\_pipeline.py}）。}
\label{tab:gates}
\small
\begin{tabularx}{\linewidth}{@{}p{1.9cm}X@{}}
\toprule
阶段 & 通过条件 \\
\midrule
1 资料搜集 & 七个类别全部覆盖；资料搜集检索 $\ge 20$ 次；近 90 天内发布的来源 $\ge 3$ 个；\code{verify-sources} 无问题 \\
2 财务抽取 & 拒收、冲突、勾稽失败均为 0，且校验晚于修改 \\
3 分析框架 & \file{framework.md} 含“外部调研清单”，“深度专题” $\ge 3$ 个 \\
4 深度调研 & 专题笔记 $\ge 3$ 份，每份 $\ge 1500$ 字且 $\ge 3$ 个独立来源；补充检索问题 $\ge 8$ 个；来源表晚于分析框架 \\
5 经营数据 & \code{verify-extra} 拒收 0 格 \\
6 估值假设 & 有参考股价及其来源；依据中引用了资料包的 \code{[S编号]} \\
7 估值运行 & 运行晚于全部输入，且带上了经营数据与资料搜集得到的可比公司表 \\
8 文稿 & 护栏（含来源）通过且核对的是当前稿；正文 $\ge 15{,}000$ 字，深度研究占比 $\ge 60\%$ \\
9 复核改写 & 初稿、修改记录与复核清单齐全；数字锁通过 \\
10 排版 & PDF 晚于文稿与脚本；25～32 页；图表 $\ge 60$ 个；已按公司与日期命名 \\
11 交付说明 & \file{delivery.md} 存在 \\
\bottomrule
\end{tabularx}
\end{table}

需要强调的是，篇幅、页数、图表数与来源数都是\textbf{数量下限}，不是质量认证；研报质量由第 \ref{sec:eval} 节的盲评兜底。数量闸门的作用是阻止模型用“写短一点”来逃避深度调研。

\subsection{先搜集、后分析}
资料搜集阶段要求覆盖七个与行业无关的类别：年报之后的最新定期报告与经营数据、最新股价与市值、卖方一致预期与目标价、可比公司及其估值倍数、行业规模与景气指标、竞争格局与份额、近三个月的重大事件。每个类别至少有一个已找到的问题，或在未找到时至少换过两个检索词。分析师读完资料包再写分析框架：3～4 条投资主线、驱动清单（年报位置 → 经营数据表 → 报告章节 → 预测参数）、外部调研清单和 3～5 个深度专题。深度专题从投资主线和资料中的分歧倒推，不预设行业套路；每个专题在阶段 4 写成一份带多个独立来源的笔记，在文稿中成为一节。

\subsection{写作规范}
\label{sec:writing}
主笔技能对“深度”提出了可检查的要求，下面几条是从人工质检中提炼的：
\begin{itemize}
  \item \textbf{结构。}年报点评约占三分之一（业绩、盈利、资产负债与现金流、年报之后的最新进展），深度研究约占三分之二（行业空间、竞争对标、核心竞争力、深度专题、预期差、估值交叉验证）。
  \item \textbf{投资主线六项要求。}每条主线必须：落到可度量的驱动；说明与一致预期差在哪个变量、差多少；落到估值参数及其对每股价值的影响（引用敏感性、情景或反向 DCF 产物）；写出可观察的证伪条件；给出检验时点；回应最强的反方论据。补不齐的降为跟踪事项。公司判断与股价判断分开写。
  \item \textbf{预测依据。}关键预测参数必须写成“依据数字 → 换算 → 取值”的链条，排版工具包读取依据表时缺项直接报错。
  \item \textbf{跟踪表。}每条核心观点给出可逐期核对的指标及“确认 / 预警 / 证伪”三档数值阈值，阈值取自产物（基准与悲观情景假设、历史区间），因而同样过得了护栏。
  \item \textbf{措辞分层。}事实直接陈述（数字来自真值表）；推论用“主要系”“反映”；观点用“我们认为”“我们预计”。不使用“事实层面”“（观点）”之类的行内标签。
\end{itemize}

\subsection{复核编辑}
复核编辑会话以“没有参与写作的研究所质控编辑”身份工作，读者设定为时间很少的买方基金经理。它先通读全文、不动笔，回答三个问题（核心结论是否一句话说得出、每条结论前是否有多环定量推导、故事线是否连贯），写在修改记录开头；再逐段问“去掉这一段读者会少知道什么”；然后做前后对账（同一指标只有一个数、结论与相邻图表一致、同一判断在摘要、正文与风险提示中方向一致）；最后整篇重读。技能说明列出了模型初稿中反复出现的机器腔写法——否定对照、自证与圆场、交代缺失、程序口吻、口语套话、数字罗列、无来源判断——但明确说明它们是判断线索而不是替换清单。需要补数字或改结论的问题写进人工复核清单 \file{review.md}，随交付物交给人。

\subsection{排版}
排版工具包 reportkit 提供读产物、标准图表与画图原语（含时间轴、流程图、文字表格等非数据图）、目录与图表目录、附录与免责声明、HTML 打印为 PDF 的能力。所有图表数据都从运行产物的 JSON 读取，不经过模型转抄。目录页码通过两遍打印回填；排版会话逐页查看渲染预览，修正留白、重叠与孤立标题。附录包含估值方法说明、关键假设汇总、全文数字的事实 / 推论 / 观点构成与统一免责声明。
